% !TeX root = ../thuthesis-example.tex

\chapter{相关工作}

本章对与本文研究相关的工作进行综述，主要包括序列推荐、生成式推荐、语义ID构建方法和多模态推荐四个方面。

\section{序列推荐}

序列推荐旨在根据用户的历史交互序列建模其动态偏好，预测下一个可能交互的物品。早期的序列推荐方法基于马尔可夫链，假设下一个交互仅依赖于最近的少量交互\cite{rendle2010factorizing}。随着深度学习的发展，基于神经网络的序列推荐方法取得了显著进展。

GRU4Rec\cite{hidasi2016session}首次将门控循环单元（GRU）引入基于会话的推荐，通过循环神经网络建模交互序列的时序依赖关系。Caser\cite{tang2018personalized}则从另一个角度出发，将用户的近期交互嵌入矩阵视为"图像"，利用卷积神经网络捕捉序列中的局部模式。

SASRec\cite{kang2018self}是序列推荐领域的里程碑工作，它将自注意力机制引入序列推荐，通过单向Transformer编码器对用户历史交互序列进行建模。SASRec使用因果注意力掩码确保每个位置只关注之前的交互，有效捕捉了长距离依赖关系。在此基础上，BERT4Rec\cite{sun2019bert4rec}借鉴BERT的掩码语言模型思想，采用双向自注意力和完形填空训练策略，在部分场景下取得了进一步提升。

值得注意的是，上述方法均基于独立的物品ID进行嵌入学习，物品之间的语义关联完全依赖于协同过滤信号通过训练隐式习得。这种方式在冷启动场景下表现不佳，且无法利用物品的内容信息。本文使用SASRec作为协同过滤嵌入的提取器，其学习到的物品嵌入编码了用户行为模式中的协同信号，为COMET融合模块提供了重要的输入。

\section{生成式推荐}

生成式推荐将推荐问题重新表述为序列生成任务。其核心思想是：首先为每个物品分配一个由多层离散编码组成的语义ID，然后训练生成模型根据用户历史交互物品的语义ID序列，自回归地生成下一个物品的语义ID。

TIGER\cite{rajput2023recommender}是生成式推荐的代表性工作。它使用RQ-VAE将物品的文本嵌入量化为多层语义ID，然后基于T5\cite{raffel2020exploring}编码器-解码器架构，将用户历史中所有物品的语义ID拼接为输入序列，自回归地生成目标物品的4层语义ID。TIGER还提出了基于前缀树（Trie）的约束束搜索，确保生成的语义ID序列一定对应某个真实物品，有效避免了无效生成。

P5\cite{geng2022recommendation}从另一个角度探索生成式推荐，将多种推荐任务（评分预测、序列推荐、解释生成等）统一为自然语言生成任务，使用单一的T5模型处理。LC-Rec\cite{zheng2024adapting}提出了对齐-生成两阶段训练范式，通过协同语义对齐提升大语言模型在推荐任务上的表现。LETTER\cite{wang2024learnable}则专注于可学习的物品标记化方法，通过引入协同过滤信号改进语义ID的构建质量。

OneRec\cite{onerec}是另一种代表性的生成式推荐模型，采用了不同于T5的自定义Transformer架构。OneRec使用交叉注意力机制将用户历史编码为键值（Key-Value）缓存，解码器通过交叉注意力查询历史上下文并自回归生成目标语义ID。其架构借鉴了大语言模型的设计，使用RMSNorm\cite{zhang2019root}、SiLU门控前馈网络\cite{shazeer2020glu}和更大的模型容量（1024维、12层）。

生成式推荐的性能很大程度上取决于语义ID的质量。如果语义ID不能有效编码物品特征，生成模型即使结构再先进也难以学习到有效的推荐模式。因此，本文聚焦于语义ID构建方法的改进。

\section{语义ID构建方法}

语义ID构建的目标是为每个物品分配一个由多层离散编码组成的标识符，使得语义相似的物品具有相近的编码。目前主流的方法基于残差量化（Residual Quantization）。

\subsection{RQ-VAE}

残差量化变分自编码器（RQ-VAE）\cite{lee2022autoregressive}是语义ID构建的基础框架。其结构包含编码器、残差向量量化器和解码器三部分。

\begin{figure}[htbp]
  \centering
  \includegraphics[width=0.8\textwidth]{RQ-VAE.png}
  \caption{RQ-VAE架构示意图\cite{lee2022autoregressive}}
  \label{fig:rq-vae}
\end{figure}

编码器将物品的高维嵌入（如4096维的文本嵌入）通过多层感知机映射到低维隐向量$\mathbf{z} \in \mathbb{R}^{d}$（如$d=32$）。残差向量量化器包含$L$层，每层维护一个码本$\mathbf{C}^{(l)} = \{\mathbf{c}_1^{(l)}, \ldots, \mathbf{c}_K^{(l)}\} \subset \mathbb{R}^{d}$（$K$通常为256）。量化过程按层级进行：
\begin{align}
  \mathbf{r}^{(0)} &= \mathbf{z}, \\
  k^{(l)} &= \arg\min_{k} \|\mathbf{r}^{(l-1)} - \mathbf{c}_k^{(l)}\|^2, \quad l = 1, \ldots, L, \\
  \mathbf{r}^{(l)} &= \mathbf{r}^{(l-1)} - \mathbf{c}_{k^{(l)}}^{(l)},
\end{align}
其中$\mathbf{r}^{(l)}$为第$l$层的残差，$k^{(l)}$为第$l$层分配的编码索引。物品的语义ID即为$(k^{(1)}, k^{(2)}, \ldots, k^{(L)})$。

解码器将量化后的表示$\hat{\mathbf{z}} = \sum_{l=1}^{L} \mathbf{c}_{k^{(l)}}^{(l)}$重建为原始嵌入。训练损失为重建损失与承诺损失之和：
\begin{equation}
  \mathcal{L}_{\text{RQ-VAE}} = \|\mathbf{x}_{\text{rec}} - \mathbf{x}\|^2 + \beta \sum_{l=1}^{L} \|\mathbf{r}^{(l-1)} - \text{sg}[\mathbf{c}_{k^{(l)}}^{(l)}]\|^2,
\end{equation}
其中$\text{sg}[\cdot]$表示停止梯度操作，$\beta$为承诺损失权重。码本通过指数移动平均（EMA）更新，并配合死码重启机制维持码本利用率。

RQ-VAE通过端到端训练优化编码质量，相比简单的聚类方法能够学习到更适合重建任务的表示。然而，RQ-VAE仅利用文本模态信息，且训练目标仅为重建文本嵌入，未考虑推荐任务中的协同过滤信号和其他模态信息，这限制了其构建的语义ID在推荐场景下的表现。

\subsection{RQ-KMeans}

RQ-KMeans\cite{onerec}是一种更简单直接的语义ID构建方法，它直接在原始嵌入空间上进行残差聚类，无需训练编码器和解码器。

\begin{figure}[htbp]
  \centering
  \includegraphics[width=0.8\textwidth]{RQ-Kmeans.png}
  \caption{RQ-KMeans架构示意图\cite{onerec}}
  \label{fig:rq-kmeans}
\end{figure}

具体而言，RQ-KMeans使用faiss库的KMeans算法对物品嵌入进行多层残差聚类。第一层对所有物品的原始嵌入$\mathbf{x} \in \mathbb{R}^{D}$（如4096维文本嵌入）进行KMeans聚类，得到$K$个聚类中心，每个物品被分配到最近的聚类中心，该聚类索引即为第一层编码$k^{(1)}$。随后计算残差$\mathbf{r}^{(1)} = \mathbf{x} - \mathbf{c}_{k^{(1)}}^{(1)}$，在残差空间上继续进行KMeans聚类得到第二层编码，如此迭代$L$层。

RQ-KMeans的优势在于方法简单、无需训练、稳定性高，但其局限性也很明显：（1）直接在高维原始空间聚类，无法通过学习优化编码质量；（2）无法引入协同过滤等额外监督信号；（3）聚类过程完全基于欧氏距离，可能无法捕捉语义相似性。尽管如此，RQ-KMeans在实践中仍能取得有竞争力的结果，常被用作基线方法。

\subsection{LETTER}

LETTER\cite{wang2024learnable}在RQ-VAE的基础上做了两项重要改进：

\begin{figure}[htbp]
  \centering
  \includegraphics[width=0.8\textwidth]{LETTER.png}
  \caption{LETTER架构示意图\cite{wang2024learnable}}
  \label{fig:letter}
\end{figure}

第一，\textbf{引入协同过滤对比损失}。LETTER使用预训练的SASRec模型提取物品的协同过滤嵌入$\mathbf{e}_{\text{cf}}$，然后在训练损失中添加InfoNCE对比损失：
\begin{equation}
  \mathcal{L}_{\text{cf}} = -\log \frac{\exp(\text{sim}(\hat{\mathbf{z}}, \mathbf{e}_{\text{cf}}) / \tau)}{\sum_{j=1}^{B} \exp(\text{sim}(\hat{\mathbf{z}}, \mathbf{e}_{\text{cf},j}) / \tau)},
\end{equation}
其中$\text{sim}(\cdot, \cdot)$为余弦相似度，$\tau=0.07$为温度参数，$B$为批大小。这一损失要求量化后的语义表示与对应的协同过滤嵌入在表示空间中对齐。为避免对比损失干扰编码器和码本的初始化，LETTER采用了预热-渐变调度：前50个训练周期不使用协同过滤对比损失，之后在50个周期内线性增加其权重至目标值。

第二，\textbf{码本多样性损失}。LETTER对每层码本的码向量进行KMeans聚类，然后在同一聚类内的码向量之间施加对比学习损失，鼓励聚类内部的码向量保持可区分性，促进码本的充分利用。

LETTER的总损失为：
\begin{equation}
  \mathcal{L}_{\text{LETTER}} = \mathcal{L}_{\text{rec}} + \mathcal{L}_{\text{vq}} + \alpha \mathcal{L}_{\text{cf}} + \gamma \mathcal{L}_{\text{div}},
\end{equation}
其中$\alpha$为协同过滤对比损失权重（经调度后为0.05），$\gamma$为多样性损失权重（0.0001）。

LETTER的局限性在于：协同过滤信号仅通过辅助损失间接影响编码过程，需要额外引入对比损失函数并仔细调节其权重；同时仍仅使用文本模态作为编码器输入，未利用图像信息。

\section{多模态推荐}

多模态信息在推荐系统中的应用日益受到关注。近年来，随着预训练多模态模型（如CLIP\cite{radford2021learning}）的发展，研究者开始探索如何有效融合文本、图像等多种模态信息以提升推荐性能。

MQL4GRec\cite{mai2025multimodal}提出了多模态量化语言的概念，旨在将不同模态的物品内容转化为统一的离散编码序列，从而实现跨域推荐中的知识迁移。该方法的核心是量化翻译器，由模态编码器和向量量化器组成。对于文本模态，MQL4GRec使用预训练语言模型提取文本嵌入，然后通过残差量化变分自编码器将其映射为离散的token序列；对于图像模态，使用预训练的视觉编码器提取图像特征，同样经过残差量化变分自编码器生成离散编码。这样，不同模态的物品内容都被转化为统一的量化语言表示。为了促进模态间的知识交互，MQL4GRec设计了量化语言对齐任务，通过对比学习显式地对齐文本和图像的量化表示，使得语义相似的文本和图像在量化语言空间中具有相近的编码。此外，该方法还设计了下一物品生成和非对称物品生成等自监督任务，使量化语言能够编码推荐相关的语义信息和先验知识。实验表明，MQL4GRec在跨域推荐任务中取得了显著的性能提升，验证了统一量化语言表示的有效性。

与MQL4GRec类似，Q-Bert4Rec\cite{qbert4rec2024}也将语义ID的思想引入多模态序列推荐。该方法首先使用预训练的文本编码器和图像编码器分别提取物品的文本嵌入和图像嵌入，然后将两种模态的嵌入进行加权融合，得到统一的多模态表示。接着，通过残差向量量化将融合后的连续嵌入映射为多层离散编码，每一层量化都在前一层的残差上进行，最终形成多层语义ID。在序列推荐阶段，采用双向Transformer架构，将用户历史交互物品的语义ID序列作为输入，通过掩码语言模型的方式预测被掩码位置的物品语义ID。实验表明，该方法在多个数据集上取得了较好的推荐性能。

上述工作表明，多模态信息的融合能够有效提升推荐性能。然而，这些方法主要关注如何在推荐模型层面融合多模态特征，或是将不同模态分别量化后进行对齐。本文则从不同的角度出发，在语义ID的编码阶段就将文本、图像和协同过滤信号通过交叉注意力机制深度融合，使生成的语义ID能够同时编码多种模态的信息，从而为下游生成式推荐模型提供更高质量的物品表示。
